\scanpage{65}
\mspage{28}
\uline{12} Récapitulation (des données de déploiement,
relatif à un ens ordonné $I$, fini disons\ldots)

\[
\left\{\begin{aligned}
\text{a) }&\text{Espaces }\Sigma_i\quad(i\in I)\\
\text{b) }&\text{Espaces }\Sigma_{i,j}\quad(i<j,\ i,j\in I)\\
\text{c) }&\text{Applications }\Sigma_{i,j}\xrightarrow{s_{ij}}\Sigma_i,\quad\Sigma_{ij}\xrightarrow{b_{ij}}\Sigma_j
\end{aligned}\right.
\]
\margin{C'est la structure envisagée dans le Scholie (p. $26/27$)
plus une application continue strictement croissante
$\Sigma\xrightarrow{\varphi}I$, telle que $\forall x\in\Sigma$, $\varphi|\Sigma_{<x}$
soit injective (\uline{donc} la relation d'ordre de \inserted{\unsure{(pour $i$ ?)}}
$\Sigma_{<x}$ soit induite par celle de $I_{<i}$).}

Conditions

a) $\Sigma_{ij}\longrightarrow\Sigma_j$ \deleted{immersion}\inserted{plongement} fermé \deleted{bordant}
\inserted{[On identifie $\Sigma_{ij}$ à un sous-espace fermé de $\Sigma_j$ par $b_{ij}$]}

b) $\Sigma_{i,k}\cap\Sigma_{j,k}=\varnothing$ si $i,j$, \deleted{$\unsure{\exists}$} non comparables
\deleted{\unsure{(on suppose}}\deleted{\unsure{$i<j$}} (NB $i,j<k$)
\margin{\uline{NB} La condition b) équivaut à celle
que dans l'ens ordonné $\Sigma=\coprod\Sigma_i$, les ensembles
$\Sigma_{<x}$ sont tot${}^t$ ordonnés.}

c) Pour $i<j<k$ posons
\[
\uline{\Sigma_{ijk}=\Sigma_{ik}\cap\Sigma_{jk}\subset\Sigma_k}.
\]
On a alors \inserted{dans $\Sigma_{jk}$ :}\deleted{$\unsure{\Sigma_{jk}}$}
\[
\uline{\Sigma_{ijk}\deleted{\unsure{\ne|\subset}}=\deleted{\unsure{\Sigma_{ik}}}b_{ij}^{-1}(\Sigma_{ij})}\deleted{\unsure{défini}}
\]
\deleted{\unsure{une}}
On a donc, pour $i<j<k$, un diagramme
commutatif hexagonal
\diagram{\begin{tikzcd}[column sep=small,row sep=large]
&\Sigma_{ik}\arrow[rr,hook,"b_{ik}"]\arrow[dl,"s_{ik}"']&&\Sigma_k&\\
\Sigma_i&&\Sigma_{ijk}\arrow[ul,hook]\arrow[rr,hook]\arrow[dl]\arrow[ur,phantom,"\text{cart}" description]\arrow[dr,phantom,"\text{cart}" description]&&\Sigma_{jk}\arrow[ul,hook,"b_{jk}"']\arrow[dl,"s_{jk}"]\\
&\unsure{\Sigma_{ij}}\arrow[ul,"s_{ij}"]\arrow[rr,hook,"b_{ij}"']&&\Sigma_j&
\end{tikzcd}}
(41)

où les flèches $\hookrightarrow$ sont des inclusions
fermées \inserted{cf plus bas} (bordantes), et les flèches $\longrightarrow$
sont des applications qui seront fibrantes et propres
dans le cas équisingulier.
\mspage{64}
